% !TEX root = main.tex
% Motivation / Introduction section for the Attack4dLLM paper.
% Notation: \method denotes our attack (the search-based diffusion red-team).
% Define once in the preamble:  \newcommand{\method}{\textsc{Search}\xspace}

\section{Motivation}
\label{sec:motivation}

\paragraph{Diffusion LLMs are a fundamentally different attack surface.}
Recent diffusion language models---\llada, \llada\,1.5, \llada\,2.0-mini,
Dream-Instruct, and \mmada-CoT---decode by iteratively denoising a
\emph{partially-masked} sequence rather than autoregressively sampling tokens
left-to-right. Two properties of this paradigm have direct safety implications.
\textbf{(i) Native infilling.} A diffusion LLM (dLLM) is trained, by
construction, to fill arbitrary $\langle\textsc{mask}\rangle$ spans conditioned on
the surrounding context. The mask token is not an out-of-distribution input but
a \emph{first-class} signal that the model is optimized to complete fluently
and confidently. \textbf{(ii) Parallel, non-causal decoding.} All masked
positions are denoised jointly across iterations, so safety filtering that
relies on \emph{prefix-conditioned} refusal (``the model has not yet committed
to a harmful direction, so it can still back out'') no longer applies: by the
time the first token is decided, the model has already been steered by the
denoising trajectory of every other position.

\paragraph{Existing red-team attacks were designed for autoregressive LLMs.}
\textsc{Pair}~\citep{chao2024pair}, \textsc{ReNeLLM}~\citep{ding2024renellm},
\textsc{AutoDAN}~\citep{liu2024autodan}, and \textsc{AIM}-style template
attacks operate on a single contiguous prompt and depend on either
gradient/log-prob signals through a causal stack or on prompt rewrites that
exploit AR refusal heuristics. None of these methods exploit the dLLM's native
mask-infill behaviour, and our experiments (Section~\ref{sec:exp}) confirm that
they transfer poorly to dLLM victims: \textsc{Pair} drops to $\le 32\%$ ASR-LLM
on HarmBench averaged across five dLLMs, and \textsc{ReNeLLM} plateaus at
$\sim 71\%$ on HarmBench but at $\le 55\%$ on JailbreakBench.

\paragraph{First-generation dLLM attacks leave headroom.}
\textsc{DIJA}~\citep{} and \textsc{PAD}~\citep{} are the only published dLLM-aware
attacks we are aware of. \textsc{DIJA} hand-crafts a small set of mask-bearing
templates and applies them uniformly; \textsc{PAD} appends paragraph-level
masks to a fixed scaffold. Both treat templating as a \emph{static} design
choice---there is no per-goal adaptation, no exploration over alternative
template structures, and no recovery mechanism when a template is rejected by
an aligned victim. Empirically this caps their ASR-LLM at
$\sim 64\%$ (\textsc{DIJA}) and $\sim 58\%$ (\textsc{PAD}) on JailbreakBench
and leaves a long tail of goals on which \emph{every} fixed template fails.

\paragraph{Our perspective: red-teaming a dLLM is a search problem.}
The space of mask-bearing templates that a dLLM will fluently complete is
enormous, and which template succeeds depends jointly on the goal, the
victim's alignment posture, and the structural prior the victim has absorbed
during pre-training (procedural lists, narrative framings, worked examples,
\ldots). No single template generalizes. Rather than hand-design a small
template family, we cast attack template generation as an
\emph{online search} over a growing pattern library:
\begin{itemize}
\itemsep0pt
  \item A \emph{pattern} is a structural schema---slot count, slot roles,
        organization style, applicable goal types---abstracted away from any
        specific surface text.
  \item For each goal, an \emph{attacker LLM} is conditioned on a selected
        pattern and instantiates a concrete mask-bearing template. The dLLM
        victim fills the masks; a \emph{reward model} scores the result.
  \item Pattern selection is governed by a \emph{UCB bandit} that trades off
        exploiting patterns with high empirical reward against exploring
        rarely-tried ones.
  \item When the bandit's best pattern fails for a hard goal, a
        \emph{scorer-guided fallback} drafts a candidate response with a
        weaker base model, redacts the harmful spans into masks, and feeds
        the redacted skeleton back to the victim---turning a single
        adversarial example into a new, goal-specific template.
  \item Templates that improve on their parent are distilled by a
        \emph{summarizer} into a new pattern schema and added to the shared
        registry, so \emph{later goals inherit the templates discovered by
        earlier goals}.
\end{itemize}

\paragraph{Why this design.}
Three observations motivated the specific choices above.
\textbf{(M1) Pattern-level abstraction beats prompt-level abstraction.}
Surface-level prompt rewrites (\textsc{ReNeLLM}-style) fail to transfer
across goals because they overfit to the lexical surface of one harmful
request. Pattern-level abstraction---``a five-slot procedural list whose
slots play (goal, premise, step, step, conclusion) roles''---transfers
because it captures the structural prior the dLLM is willing to complete
fluently, independent of which goal is being attacked. The library
therefore grows over goals rather than being thrown away after each goal.
\textbf{(M2) UCB rather than greedy.} Early in a run, every pattern has
zero visits and the empirical means are uninformative; greedy selection
collapses to whichever pattern happened to win first. UCB's $\sqrt{2\ln t/n}$
exploration term forces the bandit to retry under-visited patterns even
when one early pattern looks dominant, which we found necessary for
multi-victim transfer where the ``best'' pattern differs by victim.
\textbf{(M3) Fallback closes the long tail.} A purely template-driven
attack hits a hard ceiling on goals whose harmful realization does not
factor cleanly into the chosen pattern's slot roles. Letting a weaker
\emph{base} model (no safety alignment, but also no template structure)
draft a raw response and then \emph{re-introducing} the masks where the
harmful content lives gives the attacker a goal-conditioned template that
the bandit alone could not have produced, and is the source of the
$\sim 15$ percentage-point ASR gap between our system and the strongest
prior dLLM attack on the hardest behaviour categories.

\paragraph{Scope of contribution.}
This paper does not propose a new model architecture nor a new defence; it
provides (1) the first systematic red-team evaluation of five public dLLMs
across three benchmarks, (2) a search-based attack pipeline (\method) that
elevates ASR-LLM by $11$--$28$ points absolute over the strongest prior
dLLM attack on HarmBench / JailbreakBench / StrongReject, (3) a paragraph-mask
variant (\textsc{Dual100}) suitable for collecting safety-alignment training
data, and (4) ablations against three published defences (\textsc{A2D},
\textsc{PO}, system-prompt) showing that the search-based formulation, not
any single template family, is what drives the gain.
